\subsection{The original projective subvariety as an integral algebraic scheme}

\begin{lemma}\label{lem:original-prime-projective-kernel-integral}
\uses{graded:projective-homogeneous-quotient-immersion,thm:proj-domain-integral}
\lean{EnvelopeProperty.GradedQuotient.isIntegral_quotientKernelSubscheme}
\leanok
Let $R$ and $A$ be commutative rings, let $A$ be a graded $R$-algebra
$A=\bigoplus_{d\geq0}A_d$, and let $J$ be a homogeneous prime ideal of $A$.
Suppose that some $a\in A_d$, with $d>0$, lies outside $J$.
For the actual quotient scheme morphism
$i:\operatorname{Proj}(A/J)\to\operatorname{Proj}(A)$, the literal closed
subscheme defined by the kernel ideal sheaf of $i$ is integral.
\end{lemma}
\begin{proof}
The original quotient $A/J$ is a domain. The image of the specified $a$ in
$A/J$ is nonzero, homogeneous of the same positive degree, and belongs to
the irrelevant homogeneous ideal. Thus that irrelevant ideal is nonzero.
Apply the existing graded-domain integral-Proj theorem: its zero homogeneous
prime is a genuine relevant generic point, and its original stalks embed
in the corresponding domain localizations. Hence $\operatorname{Proj}(A/J)$
is integral. The original map $i$ is a closed immersion by the homogeneous
quotient closed-immersion theorem. A closed immersion is an isomorphism onto
its actual scheme-theoretic image, namely the subscheme of its kernel ideal
sheaf. Transport integrality along this genuine original-sheaf isomorphism.
No replacement by a reduced support, saturated ideal, or abstract projective
model is made.
\end{proof}

\begin{lemma}\label{lem:original-native-subvariety-integral}
\uses{lem:original-prime-projective-kernel-integral,def:pl-native-variety-product,projective:native-polynomial-structure-proper}
\leanok
Let $k$ be any field and let $X=V(I)\subseteq\mathbf P^N_k$ be a projective
subvariety, with its original homogeneous prime $I$ and a coordinate outside
$I$. Then the already defined native closed scheme $X^{\mathrm{nat}}$ is
integral. Its original structure map to $\operatorname{Spec}k$ makes this
same scheme a separated finite-type algebraic $k$-scheme.
\end{lemma}
\begin{proof}
Apply the preceding lemma to the original polynomial grading and original
homogeneous ideal $I$. The chosen coordinate outside $I$ has degree one,
so it supplies the required positive-degree element. The kernel subscheme
in the preceding lemma is definitionally the defined
$X^{\mathrm{nat}}$. Its original structure map is the original closed
immersion followed by the native polynomial-projective structure map.
The latter is proper, as already proved over arbitrary coefficient rings;
closed immersions are proper and composition preserves properness. Thus the
same structure map is proper, in particular separated and finite type.
Retain this actual scheme and this actual map when viewing it as an algebraic
$k$-scheme.
\end{proof}

\begin{remark}\label{rmk:original-native-integral-scheme-reuse}
These are intrinsic original-scheme inputs, over every original coefficient
field. No resolution or alteration theorem is reproved here. The author's existing polynomially projective strong-resolution theorem
provides a stronger genuine characteristic-zero resolution with integral,
polynomially projective final scheme when its actual polynomial-projective
factorization input is available. The present algebraic-scheme definition
retains the literal native scheme and original structure morphism needed
by that producer. Constructing its polynomial-projectivity input remains
separate from these three APIs. Alteration existence remains an external
interface, as requested by the repository user.
\end{remark}
